\section{Problem Formulation}

We formulate emotional preservation and prosodic realignment during low-resource language adaptation as a fully observable Markov Decision Process (MDP). Our setting starts from an expressive TTS model pretrained on a high-resource language, and considers the adaptation process to a low-resource target language. Within the IndexTTS2 framework, the environment is defined by the autoregressive Text-to-Semantic (T2S) generation process, where the agent (the language model policy $\pi_\theta$) generates a sequence of discrete semantic speech tokens token-by-token. At any step, the state $s$ encompasses the available conditioning information: the input text sequence, a speaker timbre embedding $c$, a duration constraint embedding $p$, and a target global emotion embedding $e$. The action $a$ represents the generated semantic token sequence, which is subsequently converted into a mel-spectrogram by the Semantic-to-Mel (S2M) flow-matching module. Although the policy operates on semantic tokens, the reward is computed from the decoded speech waveform.

Our objective is not to learn emotional speech generation from scratch, but to preserve and recover the emotional prior of the pretrained expressive model after low-resource adaptation. To reduce severe forgetting during transfer, we first assume a parameter-efficient adaptation stage that aligns the model to the target language while keeping most expressive knowledge intact. On top of this adapted model, the RL learner receives a multi-objective scalar reward signal $R(s, a)$ upon generating a complete sequence. This feedback is composed of four task-specific components: an Emotion Classification Reward to ensure categorical accuracy, a Global Emotion Intensity Reward measured by the distance to a neutral centroid in the Valence-Arousal-Dominance (VAD) space, a Local Emphasis Control Reward calculated from the relative pitch ($f_{pitch}$) and energy ($f_{energy}$) patterns of target emphasized words, and an ASR-Based Content Reward to preserve intelligibility and semantic correctness in the target language. The sequence-level reward can therefore be written as
\[
R = \lambda_1 R_{\text{emo-cls}} + \lambda_2 R_{\text{intensity}} + \lambda_3 R_{\text{emphasis}} + \lambda_4 R_{\text{ASR}},
\]
where the coefficients $\lambda_i$ balance emotional accuracy, global intensity, local prosody, and content preservation.

The optimization problem seeks to maximize the expected reward using Group Relative Policy Optimization (GRPO). Specifically, the objective function aims to optimize the policy parameters $\theta$ by maximizing $\mathbb{E}_{s \sim \mathcal{D}, a \sim \pi_\theta} [ R(s, a) \nabla_\theta \log \pi_\theta(a|s) ]$, subject to a Kullback-Leibler (KL) divergence constraint $\beta \text{KL}(\pi_\theta(\cdot|s) \| p_{\text{adapt}}(\cdot|s))$ to prevent catastrophic deviation from the adapted base model and to preserve its pretrained expressive prior. We operate under the assumption that the discrete semantic token space produced by the T2S module is sufficiently expressive to retain and recover subtle emotional intensities and local emphasis variations prior to acoustic reconstruction, even after transfer from a high-resource language to a low-resource language.
